\begin{helper}
Let \(\eta,\varepsilon_1\in\mathbb R\) and \(n\in\mathbb N\).  Put
\[
L=\log n,\quad
k_R=(1+\eta)\sqrt{nL},\quad
K=\noderef{TwoBiteNaturalIndependenceScale}(1+\eta,n),\quad
k=(K:\mathbb R),
\]
and
\[
t_1=\frac{\sqrt{nL}}{\log L}.
\]
Assume \(0<\eta\), \(0<\varepsilon_1\), \(0<n\), \(0<L\), \(0<\log L\),
\(k_R\le k\), and
\[
\frac1{(1+\eta)\log L}\le\varepsilon_1.
\]
Then
\[
t_1k\le\varepsilon_1 k^2.
\]
\end{helper}
\begin{proof}
Because \(0<\eta\), \(1+\eta>0\).  Since \(0<n\) and \(0<L\), we have
\(\sqrt{nL}>0\), so \(k_R>0\).  The assumption \(k_R\le k\) therefore gives
\(0<k\).

Using the definitions,
\[
t_1=\frac{\sqrt{nL}}{\log L}
=k_R\cdot\frac1{(1+\eta)\log L}.
\]
The threshold gives
\[
t_1\le \varepsilon_1 k_R.
\]
Multiplying by the nonnegative number \(k\) yields
\[
t_1k\le \varepsilon_1 k_Rk.
\]
Finally, \(k_R\le k\) and \(\varepsilon_1>0\), so
\[
\varepsilon_1 k_Rk\le \varepsilon_1 k^2.
\]
Combining the two inequalities proves the claim.
\end{proof}
